\documentclass[11pt]{amsart}
\usepackage[margin=1.1in]{geometry}
\usepackage{amsmath,amssymb,amsthm,mathtools}
\usepackage{enumitem}
\usepackage{microtype}
\usepackage{hyperref}
\usepackage[nameinlink,capitalize]{cleveref}
\hypersetup{colorlinks=true,linkcolor=blue,citecolor=blue,urlcolor=blue}

\theoremstyle{plain}
\newtheorem{theorem}{Theorem}[section]
\newtheorem{lemma}[theorem]{Lemma}
\newtheorem{proposition}[theorem]{Proposition}
\newtheorem{corollary}[theorem]{Corollary}
\theoremstyle{definition}
\newtheorem{definition}[theorem]{Definition}
\newtheorem{assumption}[theorem]{Assumption}
\theoremstyle{remark}
\newtheorem{remark}[theorem]{Remark}

\title[Local Multibubble Exclusion]{Local Multibubble and Cascade Exclusion for Type II Navier--Stokes Blowup}
\author{}
\date{}

\begin{document}

\begin{abstract}
We prove a local multibubble and cascade reduction for the conditional
Type II Navier--Stokes exclusion problem.  The setting is a positive local
Caffarelli--Kohn--Nirenberg concentration sequence at a possible singular
point.  All arguments are local in parabolic cylinders.  Compact-cylinder
failure, splitting of active CKN mass, strict same-point scale cascades,
separated active physical points, and localized gauge degeneracy are reduced by
explicit active-frame decoupling.  Under the single-core, scale-rigidity, and
terminal active-frame admissibility inputs stated below, exterior profiles and
remainders are perturbative in every selected active frame, and the remaining
branch is single-core, scale-rigid, or no longer an independent multibubble
mechanism.  Thus no local multibubble, cascade, or gauge-degenerate Type II
configuration remains under these inputs.
This paper uses the classical theory of suitable weak solutions of the
Navier--Stokes equations, including Leray \cite{Leray1934},
Caffarelli--Kohn--Nirenberg \cite{CaffarelliKohnNirenberg1982}, Scheffer
\cite{Scheffer1976}, and Lin \cite{Lin1998}.
\end{abstract}

\maketitle

\section*{Imported inputs}

This paper is a local reduction theorem.  It uses the following analytic inputs
from the preceding parts of the argument, and proves only the multibubble,
cascade, and terminal-frame decoupling steps.
\begin{enumerate}[label=\textup{(\roman*)}]
\item From Paper 1, the theorem \emph{Local single-core Type II criterion}: a
compact single-core finite-cost Type II branch with positive local CKN
concentration, nondegenerate scale-center selection, pressure reconstruction,
bounded compact-cylinder suitable package, bounded modulation, and vanishing
localized dissipation is excluded.
\item From Paper 2, the theorems \emph{Renormalized Navier--Stokes equation},
\emph{Localized pressure decomposition}, \emph{Modulation system}, and
\emph{Chart and repaired-gauge representation}: the renormalized equation,
pressure decomposition modulo functions of time, bounded modulation
coefficients, renormalized suitable structure, and stability of these structures
under subsequences.
\item From Paper 4, the theorem \emph{Compact-cylinder Caccioppoli estimate}
and the corollary \emph{Rough-core exclusion on compact windows}.
\item From Paper 5, the theorems \emph{Collapse with nonvanishing localized
core forces infinite scale-collapsing cost}, \emph{Scale-collapse cost
exclusion}, and \emph{Conditional autonomous reduced-limit theorem} wherever a
scale-rigid or scale-collapsing branch is produced.
\end{enumerate}

\section{Local setting and active frames}

\paragraph{Normalised local Navier--Stokes setting.}
All quantities are written in rescaled local variables \((x,t)\in Q_R=B_R\times(-R^2,0)\)
with viscosity \(\nu=1\) (critical scaling has absorbed \(\nu\)).  A local pair
\((u_n,p_n)\) on such a cylinder satisfies
\[
  \partial_tu_n-\Delta u_n+(u_n\cdot\nabla)u_n+\nabla p_n=0,\qquad \nabla\cdot u_n=0,
\]
and pressure is normalised through spatial averaging
\[
  (p_n)_{B_\rho}(t):=|B_\rho|^{-1}\int_{B_\rho}p_n(x,t)\,dx
\]
to factor out time-dependent constants.

For $r>0$, $x_0\in\mathbb R^3$, and $t_0\in\mathbb R$, let
\[
  Q_r(x_0,t_0):=B_r(x_0)\times(t_0-r^2,t_0),
  \qquad
  Q_r:=Q_r(0,0)=B_r\times(-r^2,0).
\]

Let $(u_n,p_n)$ be parabolic rescalings around a potential singular point with
positive local Caffarelli--Kohn--Nirenberg concentration on compact cylinders.
Concretely, for a base triple \((x_n,t_n,\lambda_n)\),
\[
  u_n(y,s)=\lambda_n u(x_n+\lambda_n y,\ t_n+\lambda_n^2 s),\qquad
  p_n(y,s)=\lambda_n^2 p(x_n+\lambda_n y,\ t_n+\lambda_n^2 s).
\]
For each $n$ and $R>0$, define the scaling-invariant local quantities
\[
  A_n(R):=R^{-1}\operatorname*{ess\,sup}_{-R^2<t<0}\int_{B_R}|u_n(x,t)|^2\,dx,
\]
\[
  E_n(R):=R^{-1}\int_{Q_R}|\nabla u_n(x,t)|^2\,dx\,dt,
\]
\[
  C_n(R):=R^{-2}\int_{Q_R}|u_n(x,t)|^3\,dx\,dt,
  \qquad
  D_n(R):=R^{-2}\int_{-R^2}^0\int_{B_R}|p_n-(p_n)_{B_R}(t)|^{3/2}\,dx\,dt,
\]
where
\[
  (p_n)_{B_R}(t):=|B_R|^{-1}\int_{B_R}p_n(x,t)\,dx.
\]

For an active frame \((x_n,\lambda_n)\), \(\lambda_n>0\), define the critical rescaling
\[
  \mathcal T_{x_n,\lambda_n}f(y):=\lambda_n f(x_n+\lambda_n y).
\]
If \(\phi\in L^3_\sigma(\mathbb R^3)\) is observed in this frame, its contribution is
\[
  \phi_n(y)=\mathcal T_{x_n,\lambda_n}\phi(y)=\lambda_n\phi(x_n+\lambda_n y).
\]
The critical Jacobian identity is exact:
\[
  \|\mathcal T_{x_n,\lambda_n}f\|_{L^3(E)}^3
  =\int_{x_n+\lambda_nE}|f(x)|^3\,dx
\]
for every measurable \(E\subset\mathbb R^3\).  In particular
\[
  \|\mathcal T_{x_n,\lambda_n}f\|_{L^3(\mathbb R^3)}^3
  =\|f\|_{L^3(\mathbb R^3)}^3,
\]
so \(\mathcal T_{x_n,\lambda_n}\) is an isometry on \(L^3(\mathbb R^3)\).

\begin{definition}[Bounded selected-window limit]
A selected-window subsequence has a bounded selected-window limit on a compact
cylinder $Q_R$ if there exists $u_\infty\in L^\infty(Q_R)$ such that, after
extracting a subsequence,
\[
  \forall\varepsilon>0\ \exists N:\ \forall n\ge N,\ \|u_n-u_\infty\|_{L^3(Q_R)}<\varepsilon.
\]
The limit is nonzero if $u_\infty\not\equiv0$ on $Q_R$.
\end{definition}

\begin{definition}[Local Type II sequence]
A positive local concentration sequence is called a local Type II sequence if every
selected-window rescaling has zero bounded selected-window limit on every compact
cylinder \(Q_R\):
\[
  \forall R>0,\qquad \lim_{n\to\infty}\|u_n\|_{L^3(Q_R)}=0.
\]
\end{definition}

\begin{definition}[Active frame]\label{def:active-frame}
Fix a concentrating time sequence \(t_n\uparrow T\) and a compact time interval
\(I\Subset(-\infty,0]\).  An active frame is a tuple
\[
  \mathfrak F^j=\bigl(x^j_n,\lambda^j_n,I,U^j_n,P^j_n\bigr),
  \qquad \lambda^j_n\downarrow0,
\]
where \(x^j_n\in\mathbb R^3\) and
\[
  U^j_n(y,s):=\lambda^j_n u(x^j_n+\lambda^j_n y,t_n+(\lambda^j_n)^2s),
  \qquad
  P^j_n(y,s):=(\lambda^j_n)^2p(x^j_n+\lambda^j_n y,t_n+(\lambda^j_n)^2s)
\]
on \(B_{2R}\times I\) for every fixed compact cylinder under consideration.
The frame is active at threshold \(\eta_*>0\) if for some fixed \(R_0<\infty\)
\[
  \liminf_{n\to\infty}\int_{Q_{R_0}}|U^j_n|^3\,dy\,ds\ge \eta_*^3.
\]
\end{definition}

\begin{definition}[Relations between active frames]\label{def:frame-relations}
Let \(\mathfrak F^i\) and \(\mathfrak F^j\) be two active frames.  After passing
to a subsequence exactly one of the following parameter regimes holds.
\begin{enumerate}[label=\textup{(\roman*)}]
\item They are \emph{same-point comparable-scale} if
\[
  0<c\le\frac{\lambda^i_n}{\lambda^j_n}\le C<\infty,\qquad
  \frac{|x^i_n-x^j_n|}{\max(\lambda^i_n,\lambda^j_n)}\le C
\]
for all large \(n\).
\item They form a \emph{same-point cascade} if, after relabeling,
\[
  \frac{\lambda^i_n}{\lambda^j_n}\to0,\qquad
  \frac{|x^i_n-x^j_n|}{\lambda^j_n}=O(1).
\]
The frame \(i\) is then the inner frame and \(j\) is an outer frame.
\item They are \emph{separated-point} if
\[
  \frac{|x^i_n-x^j_n|}{\lambda^i_n}\to\infty
  \quad\text{and}\quad
  \frac{|x^i_n-x^j_n|}{\lambda^j_n}\to\infty .
\]
\end{enumerate}
\end{definition}

\begin{definition}[Terminal frame and perturbative remainder]\label{def:terminal-frame}
Given a finite active family, a terminal frame is chosen by selecting an active
scale with minimal order:
\[
  \lambda^{j_*}_n\le C_j\lambda^j_n
  \qquad\text{for every active }j
\]
after passing to a subsequence.  Comparable frames at the same physical point
are grouped before this choice; if several grouped frames remain tied, one is
chosen by a fixed lexicographic ordering of their labels.  A remainder \(r_n\)
in a terminal frame is perturbative on \(B_{2R}\times I\) if
\[
  \|r_n\|_{L^\infty_I L^3(B_{2R})}\to0
\]
and its equation residual and pressure satisfy
\[
  \|f_n\|_{L^1_IH^{-1}(B_R)}+\|\nabla\pi^r_n\|_{L^1_IH^{-1}(B_R)}\to0
\]
for every compact \(B_R\times I\).
\end{definition}

\begin{definition}[Local multibubble/cascade branch]\label{def:local-multibubble}
A local multibubble/cascade branch is a positive-concentration local Type II
branch for which at least one of the following alternatives occurs:
\begin{enumerate}[label=\textup{(\roman*)}]
\item the compact-cylinder suitable package $A_n+E_n+C_n+D_n$ is unbounded on some
  fixed rescaled cylinder;
\item positive CKN mass splits into two or more comparable active cores;
\item a strict same-point scale cascade appears inside bounded rescaled cylinders;
\item separated physical points remain active in different local rescaled frames;
\item the localized scale or centering selection degenerates (no unique core
  selected).
\end{enumerate}
\end{definition}

\begin{definition}[Active profile package]\label{def:active-package}
An active profile package is a family
\[
  \{(\phi^j,x_n^j,\lambda_n^j)\}_{j\in J},
  \qquad \phi^j\in L^3_\sigma(\mathbb R^3),
\]
with finite or countable index set $J$, extracted along a positive local
concentration sequence.  In physical variables write
\[
  \Lambda_{\lambda,x_0}\phi(x):=\lambda^{-1}\phi\left(\frac{x-x_0}{\lambda}\right).
\]
For every finite \(J_0\subset J\) the package gives an expansion
\[
  u(\cdot,t_n)=\sum_{j\in J_0}\Lambda_{\lambda^j_n,x^j_n}\phi^j+r^{J_0}_n
\]
on the compact terminal windows under consideration.  After the finite active
subfamily has been selected, the corresponding remainder satisfies the
perturbative bounds of \cref{def:terminal-frame}.  The parameters are
considered only after grouping same-point comparable-scale profiles into
compound profiles as in \cref{def:compound-profile}.  A profile is active if
\[
  \|\phi^j\|_{L^3(\mathbb R^3)}\ge\eta_*>0.
\]
The package has finite critical mass when
\[
  \sum_{j\in J}\|\phi^j\|_{L^3(\mathbb R^3)}^3\le M_*<\infty .
\]
It captures positive local concentration if every compact terminal-window
sequence with
\[
  \liminf_{n\to\infty}\int_{Q_R}|u_n|^3>0
\]
either contains an active profile observed in that window or is absorbed into
the selected terminal frame; equivalently, no positive terminal-window
concentration remains in the perturbative remainder.
\end{definition}

\begin{definition}[Compound profile]\label{def:compound-profile}
Let \(J_k\subset J\) be a same-point comparable-scale class.  Choose one
representative \((x_n^{j_k},\lambda_n^{j_k})\).  After passing to a subsequence,
\[
  \rho_j:=\lim_{n\to\infty}\frac{\lambda^j_n}{\lambda^{j_k}_n}\in(0,\infty),
  \qquad
  z_j:=\lim_{n\to\infty}\frac{x^j_n-x^{j_k}_n}{\lambda^{j_k}_n}\in\mathbb R^3 .
\]
The compound profile in the representative frame is the \(L^3_\sigma\)-sum
\[
  \Phi^k(y):=\sum_{j\in J_k}\rho_j^{-1}
  \phi^j\left(\frac{y-z_j}{\rho_j}\right),
\]
whenever \(J_k\) is finite; finiteness follows for active classes from
\cref{lem:finite-active-count}.  If \(\|\Phi^k\|_{L^3}<\varepsilon_\star\), where
\(\varepsilon_\star\) is the critical small-data threshold, the class is
perturbative.  Otherwise \(\Phi^k\) is treated as one active object in its
representative frame.
\end{definition}

\begin{lemma}[Finite active profile count]\label{lem:finite-active-count}
Every active profile package with finite critical mass contains only finitely many
active profiles.  More precisely,
\[
  \#J\le M_*\eta_*^{-3}.
\]
\end{lemma}

\begin{proof}
For each active profile $j\in J$,
\[
  \|\phi^j\|_{L^3(\mathbb R^3)}^3\ge\eta_*^3.
\]
Let $J_0\subset J$ be finite. Summing over $J_0$ gives
\[
  \#J_0\,\eta_*^3
  \le
  \sum_{j\in J_0}\|\phi^j\|_{L^3(\mathbb R^3)}^3
  \le
  \sum_{j\in J}\|\phi^j\|_{L^3(\mathbb R^3)}^3
  \le M_*.
\]
Hence $\#J_0\le M_*\eta_*^{-3}$.  Taking the supremum over all finite
subsets $J_0$ yields the same bound for the cardinality of $J$, so $J$ is finite
and
\[
  \#J\le M_*\eta_*^{-3}.
\]
\end{proof}

\begin{lemma}[Exhaustive parameter partition]\label{lem:parameter-partition}
Let \(\mathfrak F^i,\mathfrak F^j\) be two active frames.  After passing to a
subsequence, either the pair is same-point comparable-scale, same-point cascade,
or separated-point in the sense of \cref{def:frame-relations}; if none of these
relations holds in a selected compact frame, then one of the two contributions
is locally invisible there in \(L^3\).
\end{lemma}

\begin{proof}
Pass to a subsequence so that
\[
  \frac{\lambda^i_n}{\lambda^j_n}\to \ell\in[0,\infty]
\]
in the extended half-line.  If \(0<\ell<\infty\), then either
\[
  \frac{|x^i_n-x^j_n|}{\max(\lambda^i_n,\lambda^j_n)}
\]
is bounded along a further subsequence, giving the same-point comparable-scale
case, or it tends to infinity, giving separated points because the two scales
are comparable.
\par
If \(\ell=0\), then either
\[
  \frac{|x^i_n-x^j_n|}{\lambda^j_n}
\]
is bounded, giving a same-point cascade with \(i\) inner and \(j\) outer, or it
tends to infinity.  In the latter case the \(j\)-profile is locally invisible in
the \(i\)-frame by \cref{lem:separated-vanish}.  The case \(\ell=\infty\) is the
same after interchanging \(i\) and \(j\).  These alternatives exhaust all
subsequential limits of the scale ratio and normalized center distance.
\end{proof}

\begin{theorem}[Minimal bad-configuration extraction]\label{thm:minimal-bad-extraction}
Let a positive local Type II branch admit an active profile package with finite
critical mass, perturbative remainder in the sense of
\cref{def:terminal-frame}, and concentration capture in the sense of
\cref{def:active-package}.  If the branch is not already a single active
terminal frame satisfying the single-core hypotheses, then after passing to a
subsequence one of the following occurs:
\begin{enumerate}[label=\textup{(\roman*)}]
\item compact-cylinder suitable control fails, and
\cref{lem:compactness-failure} produces a bounded-cylinder concentration core;
\item there are at least two active frames in a same-point comparable-scale
class, hence a compound profile in the sense of \cref{def:compound-profile};
\item there is a same-point cascade;
\item there are separated active physical points;
\item the scale-center selection of the selected compound or terminal frame is
degenerate.
\end{enumerate}
In every case the positive local concentration is carried by an active frame or
compound profile, not by the perturbative remainder.
\end{theorem}

\begin{proof}
If compact-cylinder suitable control fails, item \((i)\) follows from
\cref{lem:compactness-failure}.  Assume therefore that the compact-cylinder
package is bounded on the windows under consideration.
\par
By concentration capture, positive local concentration cannot remain solely in
the perturbative remainder.  Hence either one selected active frame carries the
terminal-window concentration, or another active frame remains non-negligible.
If there is only one nondegenerate selected active frame and no competing active
object, the branch is precisely the single-core regime.  Since this case is
excluded by hypothesis, either the selected scale-center map is degenerate,
giving item \((v)\), or at least two active frames remain.
\par
For two active frames, apply \cref{lem:parameter-partition}.  Same-point
comparable-scale frames give item \((ii)\), same-point scale separation gives
item \((iii)\), and separated points give item \((iv)\).  The only remaining
alternative in \cref{lem:parameter-partition} is local invisibility of one
contribution in the selected compact frame; such a contribution is absorbed
into the perturbative remainder and cannot carry the retained positive local
concentration by the capture property.  This proves the extraction statement.
\end{proof}

\section{Elementary local decoupling estimates}

\begin{lemma}[Critical rescalings that escape are locally invisible]\label{lem:escaping-l3}
Let \(\phi\in L^3(\mathbb R^3)\), let \(R<\infty\), and define
\[
  W_n(y)=\rho_n\phi(z_n+\rho_n y),
\]
where \(\rho_n>0\).  Assume that the sets
\[
  z_n+\rho_nB_R
\]
either have measure tending to zero, or escape to spatial infinity in the sense
that for every compact set \(K\subset\mathbb R^3\),
\[
  (z_n+\rho_nB_R)\cap K=\varnothing
\]
for all sufficiently large \(n\).  Then
\[
  \|W_n\|_{L^3(B_R)}\to0.
\]
\end{lemma}

\begin{proof}
By the critical change of variables,
\[
  \|W_n\|_{L^3(B_R)}^3
  =\int_{B_R}\rho_n^3|\phi(z_n+\rho_n y)|^3\,dy
  =\int_{z_n+\rho_n B_R}|\phi(x)|^3\,dx.
\]
Set
\[
  m_n:=\int_{z_n+\rho_n B_R}|\phi(x)|^3\,dx.
\]
Since \(\|W_n\|_{L^3(B_R)}^3=m_n\), it is enough to show \(m_n\to0\).
\par
Fix \(\varepsilon>0\).  We split into the two structural alternatives.

\smallskip
\noindent\textit{Case 1: shrinking measure.}
If \(|z_n+\rho_nB_R|\to0\), then by absolute continuity of the integral
\(|\phi|^3\in L^1(\mathbb R^3)\) there exists \(\delta>0\) such that
\[
  |E|<\delta\Longrightarrow \int_E|\phi|^3\,dx<\varepsilon.
\]
Choose \(N\) with \(|z_n+\rho_nB_R|<\delta\) for all \(n\ge N\).  Then
\[
  m_n=\int_{z_n+\rho_nB_R}|\phi|^3\,dx<\varepsilon,\qquad n\ge N.
\]

\smallskip
\noindent\textit{Case 2: escape to infinity.}
Assume now that for every compact set \(K\subset\mathbb R^3\),
\((z_n+\rho_nB_R)\cap K=\varnothing\) for all sufficiently large \(n\).  By
absolute continuity at infinity, choose compact \(K\subset\mathbb R^3\) such that
\[
  \int_{\mathbb R^3\setminus K}|\phi|^3<\varepsilon.
\]
For every such compact \(K\), there is \(N(K)\) with
\((z_n+\rho_nB_R)\cap K=\varnothing\) for all \(n\ge N(K)\).  For these
indices,
\[
  z_n+\rho_nB_R \subset \mathbb R^3\setminus K,
\]
and therefore, for \(n\ge N(K)\),
\[
  m_n=\int_{z_n+\rho_nB_R}|\phi|^3\,dx
  \le
  \int_{\mathbb R^3\setminus K}|\phi|^3
  <\varepsilon.
\]
In either case, for every \(\varepsilon>0\), \(m_n<\varepsilon\) for \(n\) sufficiently
large.  Hence \(m_n\to0\), and therefore
\[
  \|W_n\|_{L^3(B_R)}\to0.
\]
\end{proof}

\begin{lemma}[Separated active profiles vanish in the selected frame]\label{lem:separated-vanish}
Let \((x^p_n,\lambda^p_n)\) and \((x^q_n,\lambda^q_n)\) be two active frames
with separated physical centers in the sense that
\[
  \frac{|x^q_n-x^p_n|}{\lambda^p_n}\to\infty.
\]
If \(\phi^q\in L^3(\mathbb R^3)\), then the \(q\)-profile observed in the
\(p\)-frame converges to zero locally in \(L^3\): for every \(R<\infty\),
\[
  \left\|
  \frac{\lambda^p_n}{\lambda^q_n}
  \phi^q\left(\frac{x^p_n-x^q_n}{\lambda^q_n}
       +\frac{\lambda^p_n}{\lambda^q_n}y\right)
  \right\|_{L^3(B_R)}\to0.
\]
\end{lemma}

\begin{proof}
The displayed function has the form in \cref{lem:escaping-l3} with
\[
  z_n=\frac{x^p_n-x^q_n}{\lambda^q_n},
  \qquad
  \rho_n=\frac{\lambda^p_n}{\lambda^q_n}.
\]
The corresponding physical set in the \(q\)-profile variables is
\[
  z_n+\rho_nB_R
  =\frac{x^p_n-x^q_n}{\lambda^q_n}
   +\frac{\lambda^p_n}{\lambda^q_n}B_R.
\]
For every \(y\in B_R\),
\[
  |z_n+\rho_n y|
  \ge |z_n|-\rho_nR.
\]
Set \(E_n:=z_n+\rho_nB_R\).  Then
\[
  \operatorname{dist}(E_n,0)
  =|z_n|-\rho_nR
  =\frac{\lambda^p_n}{\lambda^q_n}\left(\frac{|x^q_n-x^p_n|}{\lambda^p_n}-R\right).
\]
Since \(\frac{|x^q_n-x^p_n|}{\lambda^p_n}\to\infty\), there exists \(N_0(R)\) such that
\[
  \frac{|x^q_n-x^p_n|}{\lambda^p_n}-R\ge1,\quad n\ge N_0(R),
\]
and hence for such \(n\), \(\operatorname{dist}(E_n,0)\ge\rho_n\).
\par
There are two cases.
\begin{enumerate}[label=\textup{(\alph*)}]
\item If \(\rho_n\to0\), then
 \[
   |z_n+\rho_nB_R|=\rho_n^3|B_R|\to0.
 \]
 Then the first alternative in \cref{lem:escaping-l3} applies.
\item Otherwise, after passing to a subsequence, \(\rho_n\ge\rho_\ast>0\).  Then
 \[
   \operatorname{dist}(z_n+\rho_nB_R,0)\ge \rho_\ast\left(\frac{|x^q_n-x^p_n|}{\lambda^p_n}-R\right)\to\infty,
 \]
 so the sets escape every fixed compact subset of \(\mathbb R^3\).  The
 second alternative in \cref{lem:escaping-l3} applies.
\end{enumerate}
In both subcases the \(L^3(B_R)\)-norm of the displayed expression converges to
zero.
\end{proof}

\begin{lemma}[Strict outer scales vanish in the innermost frame]\label{lem:outer-vanish}
Let several active profiles have the same physical center and scales
\(\lambda^1_n,\ldots,\lambda^J_n\).  Suppose \(\lambda^1_n\) is an innermost
scale and
\[
  \frac{\lambda^1_n}{\lambda^j_n}\to0
  \qquad (j\ne1).
\]
Then every outer profile \(\phi^j\in L^3(\mathbb R^3)\), viewed in the
\(\lambda^1_n\)-frame, converges to zero in \(L^3(B_R)\) for each fixed
\(R<\infty\).  Any constant ambient velocity arising from the regular part of
the outer flow is absorbed into the translation modulation.
\end{lemma}

\begin{proof}
In the innermost variables, the \(j\)-th outer contribution has the form
\[
  W^j_n(y)=\frac{\lambda^1_n}{\lambda^j_n}
  \phi^j\left(\frac{x^1_n-x^j_n}{\lambda^j_n}
  +\frac{\lambda^1_n}{\lambda^j_n}y\right).
\]
Set
\[
  \rho_n:=\frac{\lambda^1_n}{\lambda^j_n}\to0,
  \qquad
  z_n^j:=\frac{x^1_n-x^j_n}{\lambda^j_n}.
\]
Then
\[
  W^j_n(y)=\rho_n\phi^j\left(z_n^j+\rho_n y\right),
\]
and by the critical change of variables,
\[
  \|W^j_n\|_{L^3(B_R)}^3
  =\int_{z_n^j+\rho_nB_R}|\phi^j(x)|^3\,dx.
\]
Since \(\rho_n\to0\), the translated domains have
\[
  |z_n^j+\rho_nB_R|=\rho_n^3|B_R|\to0.
\]
By absolute continuity of the integral for \(\phi^j\in L^3(\mathbb R^3)\),
\[
  \|W^j_n\|_{L^3(B_R)}\to0.
\]
To treat a general \(\phi^j\), let \(\phi^j_\varepsilon\in C_c^\infty(\mathbb R^3)\) with
\(\|\phi^j-\phi^j_\varepsilon\|_{L^3}<\varepsilon\).  By critical isometry,
\[
  \|W^j_n-(W^j_n)_\varepsilon\|_{L^3(B_R)}
  \le
  \|\phi^j-\phi^j_\varepsilon\|_{L^3} < \varepsilon,
\]
where \((W^j_n)_\varepsilon\) is the same scaling of \(\phi^j_\varepsilon\).
The compactly supported term has vanishing local norm by the shrinking-domain
argument above, and then letting first \(n\to\infty\), then \(\varepsilon\downarrow0\)
implies the claimed convergence.
A constant vector field from the regular part of an outer flow contributes only
a Galilean drift in the selected frame; it is therefore incorporated into the
translation modulation and does not remain as an exterior perturbation.
\end{proof}

\begin{lemma}[Compactness failure produces another core]\label{lem:compactness-failure}
If the compact-cylinder suitable package is unbounded on a fixed rescaled
cylinder, then after passing to a subsequence the local \(C_n+D_n\) quantity is
unbounded on a comparable cylinder.  Consequently the branch contains an
additional local CKN concentration core.
\end{lemma}

\begin{proof}
Fix \(r>0\).  Apply the local Caccioppoli inequality with outer radius \(2r\).
For each index \(n\), suitable weak solutions satisfy, with \(C_r\) independent
of \(n\),
\[
  A_n(r)+E_n(r)
  \le C_r\bigl(1+C_n(2r)+C_n(2r)^{2/3}+D_n(2r)\bigr),
\]
where the pressure is normalized by subtracting spatial means
\((p_n)_{B_{2r}}(t)\).
\par
Suppose that \(A_n(r)+E_n(r)+C_n(r)+D_n(r)\) is unbounded along a subsequence.
If \(C_n(2r)+D_n(2r)\) were bounded on this subsequence, the preceding estimate
would give uniform boundedness of \(A_n(r)+E_n(r)\).  Monotonicity of the local
integrals also gives \(C_n(r)+D_n(r)\le C_n(2r)+D_n(2r)\), contradicting
unboundedness of the full package on \(Q_r\).  Therefore, after extracting a
subsequence,
\[
  C_n(2r)+D_n(2r)\to\infty .
\]
\par
Define the localized CKN concentration at variable center and radius
\[
  G_n(x,t,\rho):=\rho^{-2}\int_{Q_\rho(x,t)}|u_n|^3\,dx\,ds
  +\rho^{-2}\int_{Q_\rho(x,t)}|p_n-(p_n)_{B_\rho}(s)|^{3/2}\,dx\,ds.
\]
Set
\[
  \mathbf M_n:=\sup\{G_n(x,t,\rho):0<\rho\le 2r,\ Q_\rho(x,t)\subset Q_{4r}\}.
\]
The choice \((x,t,\rho)=(0,0,2r)\) is admissible because \(Q_{2r}(0,0)\subset Q_{4r}\), so
\[
  \mathbf M_n\ge C_n(2r)+D_n(2r)\to\infty.
\]
After extracting a further subsequence, we may assume \(\mathbf M_n\ge1\) for all \(n\).
Choose \((x_n,t_n,\rho_n)\) with \(0<\rho_n\le2r\), \(Q_{\rho_n}(x_n,t_n)\subset Q_{4r}\),
such that
\[
  G_n(x_n,t_n,\rho_n)\ge\tfrac12\,\mathbf M_n\ge\tfrac12.
\]
Rescale around \((x_n,t_n,\rho_n)\):
\[
  v_n(y,s):=\rho_nu_n(x_n+\rho_n y,t_n+\rho_n^2 s),\qquad
  q_n(y,s):=\rho_n^2p_n(x_n+\rho_n y,t_n+\rho_n^2 s).
\]
Using the critical Jacobian identity,
\[
  \int_{Q_1}|v_n|^3\,dy\,ds
  +\int_{Q_1}|q_n-(q_n)_{B_1}(s)|^{3/2}\,dy\,ds
  =G_n(x_n,t_n,\rho_n)\ge\tfrac12.
\]
Hence this new frame carries positive local critical mass, i.e. a new local CKN
concentration core.  This proves the asserted additional core extraction.
\end{proof}

\section{Terminal active-frame decoupling}

\begin{definition}[Admissible terminal active package]\label{def:admissible-terminal}
Fix a compact cylinder \(B_{2R}\times I\) and a selected active frame.  A
terminal active package is admissible on this cylinder if the following hold.
\begin{enumerate}[label=\textup{(\roman*)}]
\item The active profiles satisfy the finite-mass and active-mass conditions of
\cref{def:active-package}.
\item Every nonselected active profile is either separated from the selected
physical point in the sense of \cref{lem:separated-vanish}, or has the same
physical point and a strictly outer scale in the sense of
\cref{lem:outer-vanish}, after comparable same-point scales have been grouped.
\item The rescaled remainder \(r_n\) satisfies
\[
  \|r_n\|_{L^\infty_I L^3(B_{2R})}\to0,
\]
and there exist fields \((\pi_n^r,f_n)\) such that on \(B_{2R}\times I\),
\[
  \partial_t r_n-\Delta r_n+a_n\bigl(r_n+y\cdot\nabla r_n\bigr)+b_n\cdot\nabla r_n
  +\nabla\pi_n^r=f_n,\qquad
  \nabla\cdot r_n=0,
\]
with
\[
  \|f_n\|_{L^1_IH^{-1}(B_R)}\to0,
\]
\[
  \|\nabla\pi_n^r\|_{L^1_IH^{-1}(B_R)}\to0.
\]
\item The selected velocity satisfies
\[
  \sup_n\|U_n\|_{L^\infty_I L^3(B_{2R})}<\infty .
\]
\item The selected modulation coefficients satisfy
\[
  \sup_n\bigl(\|a_n\|_{L^\infty(I)}+\|b_n\|_{L^\infty(I)}\bigr)<\infty .
\]
\item Local pressure decomposition is stable under the splitting: if tensor
coefficients tend to zero in \(L^1_IL^{3/2}(B_R)\), then the corresponding
pressure gradients tend to zero in \(L^1_IH^{-1}(B_R)\), after subtracting
spatial means.
\item Cutoff commutators generated by localizing to \(B_R\) are controlled by
the \(L^\infty_I L^3\), \(L^1_I L^{3/2}\), and pressure-gradient norms appearing
above.  In particular, for any compact \(\chi_R\in C_c^\infty(B_{2R})\) with
\(\chi_R\equiv1\) on \(B_R\), the commutators satisfy
\[
  \|[\Delta,\chi_R]u\|_{H^{-1}(B_R)}
  +\|[\nabla\chi_R,\,\nabla\cdot](u\otimes u)\|_{H^{-1}(B_R)}
  \le C_R\|u\|_{L^3(B_{2R})}^2.
\]
uniformly for \(u\in L^3(B_{2R})\) and the same type of estimate for
\(S_n\)-quadratic commutator terms.
\end{enumerate}
\end{definition}

\begin{theorem}[Terminal profile-evolution decoupling]\label{thm:terminal-decoupling}
Consider an admissible terminal active package on \(B_{2R}\times I\).  Let
\(U_n\) be the selected-frame velocity and let \(S_n\) be the sum, in that
frame, of all nonselected active profiles together with the rescaled remainder.
Then on \(B_R\times I\):
\begin{enumerate}[label=\textup{(\roman*)}]
\item \(\nabla\cdot S_n=0\) distributionally;
\item \(S_n\to0\) in \(L^\infty_I L^3(B_{2R})\);
\item \(U_n\otimes S_n+S_n\otimes U_n+S_n\otimes S_n\to0\) in
\(L^1_I L^{3/2}(B_R)\);
\item
\[
  \begin{aligned}
  &\exists\,\pi_n,\ F_n:\quad
  \partial_t S_n-\Delta S_n
  +a_n\bigl(y\cdot\nabla S_n + S_n\bigr)
  +b_n\cdot\nabla S_n+\nabla\pi_n =F_n,\\
  &\|F_n\|_{L^1_IH^{-1}(B_R)}\to0,\qquad
  \nabla\cdot S_n=0\quad\text{in }\mathcal D'(B_R\times I).
  \end{aligned}
\]
\[
  \|\nabla\pi_n\|_{L^1_IH^{-1}(B_R)}\to0.
\]
\item every pressure source containing at least one factor of \(S_n\) has
vanishing pressure gradient in \(L^1_IH^{-1}(B_R)\);
\item localization errors introduced by compact cutoffs vanish in
\(L^1_IH^{-1}(B_R)\);
\item the effective modulation parameters of the selected frame remain bounded;
\item after discarding \(S_n\), the selected frame remains a positive finite-mass
local Type II branch with pressure reconstruction and the compact-cylinder
bounds required by the single-core analysis.
\end{enumerate}
\end{theorem}

\begin{proof}
Each profile in the package is divergence-free, and the rescaled remainder is
divergence-free.  Since divergence commutes with the Navier--Stokes critical
rescalings, finite sums remain divergence-free in distributions.  This proves
\((i)\).

We prove \((ii)\).  By admissibility, every nonselected active profile is either
separated from the selected physical point or is a same-point strictly outer
scale.  The separated profiles vanish locally by \cref{lem:separated-vanish},
and the same-point outer profiles vanish locally by \cref{lem:outer-vanish}.
The remainder tends to zero in \(L^\infty_I L^3(B_{2R})\) by
\cref{def:admissible-terminal}.  The number of active profiles is finite by
\cref{lem:finite-active-count}.  Since every summand in \(S_n\)
tends to zero in \(L^\infty_I L^3(B_{2R})\), the triangle inequality gives
\[
  \|S_n\|_{L^\infty_I L^3(B_{2R})}
  \le
  \sum_{j\ne j_*}\|S^j_n\|_{L^\infty_I L^3(B_{2R})}
  +\|r_n\|_{L^\infty_I L^3(B_{2R})}\to0.
\]
This proves \((ii)\).

For \((iii)\), admissibility gives
\[
  \sup_n\|U_n\|_{L^\infty_I L^3(B_{2R})}<\infty.
\]
Hence, by H\"older's inequality,
\[
  \|U_n\otimes S_n\|_{L^1_IL^{3/2}(B_R)}
  \le |I|\|U_n\|_{L^\infty_IL^3(B_R)}
       \|S_n\|_{L^\infty_IL^3(B_R)}\to0,
\]
and the same estimate applies to \(S_n\otimes U_n\).  Finally,
\[
  \|S_n\otimes S_n\|_{L^1_IL^{3/2}(B_R)}
  \le |I|\|S_n\|_{L^\infty_IL^3(B_R)}^2\to0.
\]

For \((iv)\), write the represented Navier--Stokes equation for the full
profile sum \(U_n+S_n\) and subtract the represented equation for the selected
component \(U_n\).  The remaining equation for \(S_n\) has the displayed linear
left-hand side
\[
  \partial_t S_n-\Delta S_n
  +a_n\bigl(y\cdot\nabla S_n+S_n\bigr)+b_n\cdot\nabla S_n.
\]
Collect all right-hand terms into
\[
  F_n:= -\nabla\cdot(U_n\otimes S_n+S_n\otimes U_n+S_n\otimes S_n)
       +f_n+E_n^{\mathrm{cut}}+E_n^{\mathrm{mod}},
\]
where \(f_n\) is the remainder forcing from item \((iii)\) of admissibility, and
\(E_n^{\mathrm{cut}},E_n^{\mathrm{mod}}\) are localization and modulation errors.
The nonlinear stresses tend to zero in \(L^1_IL^{3/2}(B_R)\), hence in
\(L^1_IH^{-1}(B_R)\), because \(L^{3/2}(B_R)\hookrightarrow H^{-1}(B_R)\) on a
bounded three-dimensional ball.  The residual term vanishes by
\cref{def:admissible-terminal}.  The cutoff terms vanish by the localization
part of admissibility, recorded explicitly in \((vi)\).  This proves the
displayed convergence in \(L^1_IH^{-1}(B_R)\).

For \((v)\), decompose the pressure source of \(U_n+S_n\) into the pure selected
source, the pure exterior source, and the two mixed sources.  The mixed sources
are
\[
  \partial_i\partial_j(U_{n,i}S_{n,j}+S_{n,i}U_{n,j})
  \quad\text{and}\quad
  \partial_i\partial_j(S_{n,i}S_{n,j}).
\]
By \((iii)\), their tensor coefficients tend to zero in
\(L^1_IL^{3/2}(B_R)\).  The pressure-stability part of admissibility,
equivalently local Calderon--Zygmund estimates plus harmonic pressure control
modulo spatial means, gives convergence of the corresponding pressure gradients
to zero in \(L^1_IH^{-1}(B_R)\).  Pure exterior sources are generated by
profiles that are locally invisible in the selected frame; applying the same
pressure stability after \cref{lem:separated-vanish,lem:outer-vanish} gives
their vanishing pressure gradients.

For \((vi)\), let \(\chi_R\) be a cutoff equal to one on \(B_R\) and supported
in \(B_{2R}\).  The commutators are finite sums of terms of the following
types:
\[
  (\nabla\chi_R)\cdot\nabla S_n,
  \qquad
  (\Delta\chi_R)S_n,
  \qquad
  (\nabla\chi_R)\cdot (U_n\otimes S_n+S_n\otimes U_n+S_n\otimes S_n),
\]
plus analogous pressure cutoff terms.  Their supports lie in the fixed annulus
\(B_{2R}\setminus B_R\).  By the localization part of admissibility, these
commutators are controlled by the norms already shown to vanish in \((ii)\),
\((iii)\), and \((v)\).  Hence all localization errors vanish in
\(L^1_IH^{-1}(B_R)\).

For \((vii)\), the only exterior contribution that can remain nonzero at first
order in a compact selected frame is a constant ambient velocity.  This has
already been absorbed into the translation modulation.  The remaining exterior
profiles vanish locally, so their contribution to the modulation equations is
small in the same local \(L^3\) topology used above.  The selected modulation
coefficients are bounded by admissibility; hence the effective parameters remain
bounded.

For \((viii)\), the positive CKN mass of the selected profile is unchanged by
subtracting a term that tends to zero in local \(L^3\) and whose
mixed pressure sources vanish by \((v)\).  The pressure reconstruction is stable
under the same decomposition by the pressure-stability condition in
\cref{def:admissible-terminal}.  The compact-cylinder bounds pass to the
selected component because the full package is bounded and the discarded terms
vanish on compact cylinders.  Thus the selected frame remains a positive
finite-mass local Type II branch with the analytic inputs needed for the
single-core analysis.
\end{proof}

\section{Same-point and separated-point reductions}

\begin{theorem}[Same-point cascade reduction]\label{thm:same-point-reduction}
Every same-point multibubble or strict same-point cascade reduces, in the
selected innermost or compound frame, to one of the following: the single-core
criterion, scale-rigidity, or a perturbative terminal branch covered by
\cref{thm:terminal-decoupling}, provided the terminal package generated in the
selected innermost frame is admissible in the sense of
\cref{def:admissible-terminal}.
\end{theorem}

\begin{proof}
\noindent\textit{Step 1: scale-comparability classes.}
Split active same-point profiles into equivalence classes by
\[
  i\sim j \iff \exists c,C\in(0,\infty):
  c\le \frac{\lambda^i_n}{\lambda^j_n}\le C \quad \text{for all }n
\]
after passing to a subsequence.
Thus within each class there are constants \(c_{ij},C_{ij}\) with
\(c_{ij}\le\lambda^i_n/\lambda^j_n\le C_{ij}\) for all \(n\), so one may define
a single common rescaling sequence (choose \(\lambda_n^{(k)}=\lambda_n^{j_k}\)) and write
all class-\(k\) profiles in the same coordinates.
Within each class, finite sums of divergence-free profiles remain divergence-free.

\smallskip
\noindent\textit{Step 2: compound profile reduction inside each class.}
For each class set
\[
  \Phi^k_n:=\sum_{j\in J_k}\lambda^j_n\phi^j(x_n^j+\lambda^j_n y).
\]
If the \(L^3\)-size of this rescaled sum satisfies
\[
  \left\|\sum_{j\in J_k}\phi^j\right\|_{L^3(\mathbb R^3)}<\varepsilon_\star,
\]
with \(\varepsilon_\star\) the perturbative small-data threshold, then standard local
small-data theory for Navier--Stokes in the critical space removes this whole class
from the active list.

If the class is above threshold, it is retained as a single compound active core;
no longer than one active core is needed from that class for local selection.

\smallskip
\noindent\textit{Step 3: local alternatives.}
If the retained compound core is unique and its scale-centering Jacobian is
nondegenerate, we obtain the single selected core alternative.
If this Jacobian is degenerate, the localized choice of active center or scale is
not unique, which is the gauge-degenerate alternative in
\cref{def:local-multibubble}.

\smallskip
\noindent\textit{Step 4: strict same-point scale separation.}
Assume now the branch is not reduced to comparable scales and choose the innermost
active scale \(\lambda^1_n\) in that physical point.  By
\cref{lem:outer-vanish}, every profile with larger scale vanishes in the
\(\lambda^1_n\)-frame in local \(L^3\), up to constant drift terms.
These constant drifts are absorbed into translation modulation by admissibility.
Let \(S_n\) be the sum of all such outer profiles plus the rescaled remainder
\(r_n\).  Then, by Step 2 and \cref{lem:outer-vanish},
\[
  S_n\to0\quad\text{in }L^\infty_I L^3(B_{2R}).
\]

Subtracting the equation for the innermost selected profile from the full-frame
equation gives a perturbation equation for the selected profile whose forcing
belongs to \(L^1_I H^{-1}(B_R)\) and vanishes.  This is the setting of
\cref{thm:terminal-decoupling}.  Therefore one of three outcomes holds:
\begin{enumerate}[label=\textup{(\alph*)}]
\item \textbf{Single-core branch:}
  after discarding \(S_n\), the branch is governed by one active component, so it
  lands in the compact single-core branch of \cref{ass:single-core};
\item \textbf{Scale-rigid branch:}
  the relative Jacobian between comparable scales becomes rigid in the sense of
  \cref{ass:scale-rigidity}, so \cref{ass:scale-rigidity} excludes it;
\item \textbf{Purely perturbative branch:}
  \(S_n\to0\) and the selected profile satisfies a vanishing forcing perturbation
  in \(L^1_IH^{-1}(B_R)\).  This branch carries no independent same-point
  cascade mechanism; after discarding \(S_n\), the selected component is passed
  to the single-core or scale-rigidity analysis.
\end{enumerate}

These are the strict same-point scale-separation alternatives, so the
same-point reduction is exhausted.
\end{proof}


\begin{theorem}[Separated-point decoupling]\label{thm:separated-point-reduction}
Separated active physical points are locally invisible in each other's active
rescaled frames.  Consequently every separated-point multibubble reduces in each
selected active frame to a single-core, scale-rigid, or perturbative terminal
branch, provided the terminal package in that selected active frame is
admissible in the sense of \cref{def:admissible-terminal}.
\end{theorem}

\begin{proof}
Fix an active point \(p\) with frame \((x^p_n,\lambda^p_n)\).  Let \(q\ne p\)
be another active point.  By separation,
\[
  \frac{|x^q_n-x^p_n|}{\lambda^p_n}\to\infty.
\]
Thus \cref{lem:separated-vanish} gives local \(L^3\) convergence to zero of the
\(q\)-profile in the \(p\)-frame on every compact ball.  The same conclusion
holds for each profile centered at any physical point different from \(p\).
Since only finitely many active profiles remain after the critical mass
threshold is imposed, the sum of all profiles outside the \(p\)-frame tends to
zero locally in \(L^3\).

The pure pressure generated by a separated profile is locally invisible by the
same change of variables whenever its profile pressure belongs to
\(L^{3/2}_{\mathrm{loc}}\) with the standard Calderon--Zygmund control.  Mixed
pressure terms contain at least one exterior velocity factor, hence vanish in
\(L^1H^{-1}_{\mathrm{loc}}\) by the pressure part of
\cref{thm:terminal-decoupling}.  Localization and modulation commutators are
localized to \(B_{2R}\setminus B_R\), so their \(H^{-1}\)-norm is controlled by
the \(L^3\)-smallness already proved in \cref{lem:separated-vanish,lem:outer-vanish}
and by admissibility item \((\text{vii})\).

Therefore, in the selected \(p\)-frame, all other physical points contribute
only a perturbative forcing tending to zero in \(L^1H^{-1}\) on compact
cylinders.  The selected frame remains a positive finite-mass branch with the
same compact-cylinder and pressure reconstruction properties.  Hence the branch
reduces to the single-core case, the scale-rigid case, or the perturbative
terminal branch.  Since \(p\) was arbitrary, the same reduction holds in every
active frame.
\end{proof}

\section{Multibubble and cascade exclusion}

The assumptions in this section isolate the analytic inputs needed for the
multibubble reduction: the single-core exclusion, the scale-rigidity reduction,
and terminal active-frame admissibility.  They are stated explicitly so that the
theorem below uses no unstated compactness, pressure, or profile-completeness
conclusion.

\begin{assumption}[Local single-core exclusion]\label{ass:single-core}
Every compact single-core branch with positive local CKN concentration,
nondegenerate scale-center selection, pressure reconstruction, bounded
compact-cylinder suitable package, bounded modulation, and vanishing localized
dissipation is excluded from the local Type II regime.
\end{assumption}

\begin{assumption}[Local scale-rigidity exclusion]\label{ass:scale-rigidity}
Every scale-rigid local branch produced by the same-point or separated-point
reductions, including a branch with degenerate relative scale-center Jacobian,
either has a nonzero bounded selected-window limit or reduces to the single-core
branch in \cref{ass:single-core}.  Consequently no such branch is compatible
with the local Type II condition.
\end{assumption}

\begin{assumption}[Terminal admissibility in active frames]\label{ass:terminal-admissibility}
Every terminal package produced by the same-point and separated-point reductions
is admissible on each compact selected cylinder in the sense of
\cref{def:admissible-terminal}.
\end{assumption}

\begin{theorem}[Local multibubble and cascade exclusion]\label{thm:multi}
Assume the local single-core and scale-rigidity exclusions
\cref{ass:single-core,ass:scale-rigidity}, and assume terminal admissibility in
active frames, \cref{ass:terminal-admissibility}.  Let a positive local Type II
branch admit an active package satisfying the hypotheses of
\cref{thm:minimal-bad-extraction}.  Then no local multibubble, cascade, or
gauge-degenerate Type II case can occur in this class.
\end{theorem}

\begin{proof}
Let a local Type II branch be in the multibubble/cascade class of
\cref{def:local-multibubble}.  Apply \cref{thm:minimal-bad-extraction}.  The
positive local concentration is carried by an active frame or compound profile,
and one of the alternatives listed there occurs.
\par
\textit{Case 1: compactness-failure reduction.}
Compactness failure yields, after changing radius by a fixed factor and
rescaling as in \cref{lem:compactness-failure}, a bounded-cylinder concentration
branch.  Relabel this branch as an active core.  After passing to a subsequence,
the new core and every previously retained active core have one of three
relations: comparable scale at the same physical point, strictly separated scale
at the same physical point, or separated physical center.  Comparable same-point
cores are grouped into compound profiles.  A perturbative compound profile is
removed by small-data stability; an active compound profile is treated as one
selected core in its active frame.  If the scale-center selection inside such a
compound core is degenerate, the branch is gauge-degenerate, giving the
non-uniqueness alternative in \cref{def:local-multibubble} rather than a new
independent core dynamics.
\par
\textit{Case 1a: gauge degeneracy.}
If the scale-center selection of the selected compound or terminal frame is
degenerate, then the branch is in the gauge-degenerate alternative of
\cref{def:local-multibubble}.  By the imported repaired-gauge representation
from Paper 2, nondegeneracy is exactly the hypothesis needed to enter the
single-core branch.  Persistent degeneracy is therefore routed to the
scale-rigidity alternative in \cref{ass:scale-rigidity}; if the degeneracy
disappears after passing to a subsequence, the branch enters the single-core or
terminal-frame alternatives below.
\par
\textit{Case 1b: same-point comparable scales.}
Same-point comparable-scale active frames are grouped into the compound profile
of \cref{def:compound-profile}.  If the compound profile is below the critical
small-data threshold, it is part of the perturbative remainder.  If it is above
threshold and its scale-center selection is nondegenerate, it is one selected
active core and enters the single-core branch.  If the corresponding
scale-center selection is degenerate, it is covered by Case 1a.  Thus
comparable same-point profiles do not produce an additional multibubble
mechanism.
\par
\textit{Case 2: same-point cascading.}
Strict same-point scale cascades are handled by
\cref{thm:same-point-reduction}.  In the innermost frame, all outer scales are
locally invisible after constant drifts are absorbed.  The branch therefore
reduces to one of the three terminal alternatives
\[
  \text{(A) single-core},\quad \text{(B) scale-rigid},\quad \text{(C) terminal perturbative}.
\]
\par
\textit{Case 3: separated-point interactions.}
For separated physical points, \cref{thm:separated-point-reduction} gives the same
three terminal alternatives as in Case 2:
\[
  \text{(A) single-core},\quad \text{(B) scale-rigid},\quad \text{(C) terminal perturbative}.
\]

The single-core alternative is excluded by \cref{ass:single-core}.  The
scale-rigid alternative is excluded by \cref{ass:scale-rigidity}, which packages
the Paper 5 scale-collapse and autonomous-limit exclusions in this setting.  The
terminal perturbative alternative cannot sustain an independent multibubble
obstruction because
\cref{thm:terminal-decoupling} removes all exterior profiles and remainders from
the selected active equation in \(L^1H^{-1}_{\mathrm{loc}}\), preserves the positive
finite-mass selected branch, and leaves only the single-core or scale-rigid
possibilities already excluded.  Consequently every alternative in
\cref{def:local-multibubble} is impossible in the local Type II regime.
\end{proof}

\begin{corollary}[Removal of the multibubble alternative]\label{cor:multi-removal}
In the local Type II analysis, every failure of the single nondegenerate compact
core alternative is excluded by \cref{thm:multi} under the hypotheses stated
there.
\end{corollary}

\begin{proof}
By \cref{def:local-multibubble}, failure of the single nondegenerate compact
core alternative means that at least one of the following occurs: the
compact-cylinder suitable package is unbounded; positive local CKN mass splits
between active cores; a strict same-point scale cascade occurs; separated
physical points remain active in distinct local frames; or the localized scale
or center selection is degenerate.  The compactness-failure case produces an
additional core by \cref{lem:compactness-failure}.  Same-point and separated
cases are reduced by \cref{thm:same-point-reduction,thm:separated-point-reduction}.
The final exclusion is \cref{thm:multi}.  Hence no failure of the single
nondegenerate compact core alternative remains.
\end{proof}

\bibliographystyle{alpha}
\bibliography{paper3_critical_profile_decompositions_radiation_multibubbles}

\end{document}
